\documentclass[11pt, a4paper]{article}

% ==================== PAQUETES ====================
\usepackage[utf8]{inputenc}
\usepackage[T1]{fontenc}
\usepackage[spanish]{babel}
\usepackage[margin=2.3cm]{geometry}
\usepackage{amsmath, amssymb}
\usepackage{array, booktabs, tabularx}
\usepackage{xcolor}
\usepackage{graphicx}
\usepackage{tikz}
\usepackage{hyperref}
\usepackage{fancyhdr}
\usepackage{enumitem}
\usepackage{titlesec}
\usepackage{tcolorbox}
\usepackage{longtable}

\usetikzlibrary{shapes, arrows, positioning, calc, trees, fit}
\tcbuselibrary{skins, breakable}

% ==================== COLORES ====================
\definecolor{policiaAzul}{HTML}{0A3B66}
\definecolor{policiaAzulClaro}{HTML}{1E5A8E}
\definecolor{policiaDorado}{HTML}{C9A227}
\definecolor{grisClaro}{HTML}{F4F6F8}
\definecolor{grisBorde}{HTML}{DCE2E8}
\definecolor{acentoSuave}{HTML}{EAF2FA}
\definecolor{exitoVerde}{HTML}{1E7E34}
\definecolor{exitoSuave}{HTML}{E9F7EF}
\definecolor{alertaRojo}{HTML}{C0392B}
\definecolor{alertaSuave}{HTML}{FDECEA}
\definecolor{causa}{HTML}{A93226}
\definecolor{consecuencia}{HTML}{7D3C98}

% ==================== ESTILO ====================
\hypersetup{
  colorlinks=true,
  linkcolor=policiaAzul,
  urlcolor=policiaAzulClaro,
  pdftitle={Guia Demuestro mi aprendizaje U2 - Matematica aplicada},
  pdfauthor={Jose Javier Mercado Reyes - ESJIM}
}

\titleformat{\section}[hang]
  {\Large\bfseries\color{policiaAzul}}
  {\thesection.}{0.6em}{}
  [\vspace{4pt}{\color{policiaDorado}\titlerule[1.2pt]}]

\titlespacing*{\section}{0pt}{1.8em}{0.9em}

\titleformat{\subsection}[hang]
  {\large\bfseries\color{policiaAzulClaro}}
  {\thesubsection.}{0.5em}{}

\titlespacing*{\subsection}{0pt}{1.2em}{0.5em}

% Encabezado/pie
\pagestyle{fancy}
\fancyhf{}
\fancyhead[L]{\footnotesize\color{policiaAzul}\textbf{ESJIM} $\cdot$ Matem\'atica aplicada a la profesi\'on}
\fancyhead[R]{\footnotesize\color{policiaAzul}U2 $\cdot$ Demuestro mi aprendizaje}
\fancyfoot[C]{\footnotesize\color{policiaAzul}P\'agina \thepage\ de \pageref*{LastPage}}
\renewcommand{\headrulewidth}{0.4pt}
\renewcommand{\headrule}{\color{policiaDorado}\hrule width\headwidth height0.8pt}
\usepackage{lastpage}

% Cajas tem\'aticas
\newtcolorbox{cajaInfo}[1][]{
  colback=acentoSuave, colframe=policiaAzulClaro,
  fonttitle=\bfseries\color{white}, title=#1,
  arc=2mm, left=10pt, right=10pt, top=6pt, bottom=6pt,
  breakable
}
\newtcolorbox{cajaExito}{
  colback=exitoSuave, colframe=exitoVerde, coltitle=exitoVerde,
  arc=2mm, left=10pt, right=10pt, top=6pt, bottom=6pt,
  leftrule=4pt, rightrule=0pt, toprule=0pt, bottomrule=0pt,
  breakable
}
\newtcolorbox{cajaAlerta}{
  colback=alertaSuave, colframe=alertaRojo, coltitle=alertaRojo,
  arc=2mm, left=10pt, right=10pt, top=6pt, bottom=6pt,
  leftrule=4pt, rightrule=0pt, toprule=0pt, bottomrule=0pt,
  breakable
}
\newtcolorbox{cajaPaso}[1][]{
  colback=white, colframe=policiaDorado, coltitle=policiaAzul,
  fonttitle=\bfseries, title=#1,
  arc=2mm, left=10pt, right=10pt, top=6pt, bottom=6pt,
  leftrule=4pt, rightrule=0.4pt, toprule=0.4pt, bottomrule=0.4pt,
  breakable
}

% ==================== DOCUMENTO ====================
\begin{document}

% ---- PORTADA ----
\begin{titlepage}
\centering
\vspace*{1.5cm}

\begin{tcolorbox}[colback=policiaAzul, colframe=policiaDorado, boxrule=2pt, arc=4mm, left=20pt, right=20pt, top=20pt, bottom=20pt]
\color{white}
{\footnotesize\textbf{ESCUELA DE SUBOFICIALES Y NIVEL EJECUTIVO}}\\[2pt]
{\footnotesize ``GONZALO JIM\'ENEZ DE QUESADA''}\\[10pt]
\color{policiaDorado}
\rule{0.4\textwidth}{0.6pt}\\[10pt]
\color{white}
{\LARGE\textbf{Gu\'ia de trabajo}}\\[6pt]
{\Large \textit{Demuestro mi aprendizaje}}\\[10pt]
\color{policiaDorado}
{\large\textbf{Unidad 2}}\\[3pt]
\color{white}
{\large Operaciones B\'asicas para la Gesti\'on Efectiva}\\
{\large en la Administraci\'on Policial}\\[14pt]
\color{policiaDorado}\rule{0.4\textwidth}{0.6pt}\\[10pt]
\color{white}
{\normalsize Matem\'atica aplicada a la profesi\'on}
\end{tcolorbox}

\vspace{1.2cm}

\begin{tabular}{@{}p{0.25\textwidth}p{0.65\textwidth}@{}}
\textbf{\color{policiaAzul}Docente:} & Jos\'e Javier Mercado Reyes\\[5pt]
\textbf{\color{policiaAzul}Asignatura:} & 338-ADMI-MATE-2 Matem\'atica aplicada a la profesi\'on\\[5pt]
\textbf{\color{policiaAzul}Secciones:} & 008227 $\cdot$ 008228 $\cdot$ 008229 $\cdot$ 008230\\[5pt]
\textbf{\color{policiaAzul}Unidad:} & 2 --- Operaciones B\'asicas para la Gesti\'on Efectiva\\[5pt]
\textbf{\color{policiaAzul}Actividad:} & Demuestro mi aprendizaje\\[5pt]
\textbf{\color{policiaAzul}Apertura:} & Lunes 16 de marzo de 2026, 00:01\\[5pt]
\textbf{\color{policiaAzul}Cierre:} & Jueves 14 de mayo de 2026, 23:59\\[5pt]
\textbf{\color{policiaAzul}Peso:} & 5.00 puntos (t\'ecnica sumativa de la unidad)
\end{tabular}

\vfill

\begin{cajaInfo}[\textbf{Prop\'osito de esta gu\'ia}]
Este documento contiene todo lo que usted necesita para desarrollar la actividad ``Demuestro mi aprendizaje'' de la Unidad 2. Incluye un \textbf{ejercicio resuelto de programaci\'on lineal paso a paso}, una \textbf{plantilla de \'arbol de problemas}, una \textbf{plantilla de an\'alisis costo-beneficio}, la r\'ubrica de evaluaci\'on y las referencias bibliogr\'aficas. Use el ejercicio de ejemplo como gu\'ia para construir su propio caso.
\end{cajaInfo}

\vspace{0.5cm}
{\footnotesize \color{policiaAzul} Abril 2026}

\end{titlepage}

\tableofcontents
\newpage

% ==================== 1. OBJETIVO ====================
\section{Objetivo y competencias}

\subsection{Resultado de aprendizaje}
\begin{cajaInfo}[\textbf{RA-U2}]
Realiza c\'alculos y operaciones matem\'aticas fundamentales (porcentajes, proporciones, manejo de datos cuantitativos) en el contexto policial para tomar decisiones informadas que fortalezcan la eficiencia operativa en su labor profesional.
\end{cajaInfo}

\subsection{Competencias espec\'ificas}
Al completar esta actividad el estudiante estar\'a en capacidad de:
\begin{itemize}[leftmargin=1.5em]
\item \textbf{Estructurar} un problema real del servicio policial identificando con claridad sus causas y consecuencias.
\item \textbf{Formular} una alternativa de soluci\'on con variables medibles.
\item \textbf{Evaluar} la viabilidad de esa alternativa mediante an\'alisis costo-beneficio, ya sea con valores num\'ericos o con argumentaci\'on anal\'itica.
\item \textbf{Justificar} una decisi\'on operativa de manera que pueda ser auditada por el mando.
\end{itemize}

% ==================== 2. COMPONENTES ====================
\section{Componentes obligatorios de la entrega}

La entrega debe contener como m\'inimo los dos elementos siguientes:

\begin{cajaPaso}[\textbf{Componente 1 --- \'Arbol de problemas}]
Represente gr\'aficamente una \textbf{situaci\'on cotidiana del servicio policial} (por ejemplo: distribuci\'on del combustible mensual, rotaci\'on de turnos, elecci\'on entre mantenimiento preventivo y correctivo, instalaci\'on de c\'amaras adicionales, etc.). El \'arbol debe mostrar:
\begin{itemize}[leftmargin=1.5em]
\item \textbf{Problema central} (tronco): redactado en una frase clara, medible y ubicada en el tiempo y el territorio.
\item \textbf{Causas} (ra\'ices): m\'inimo 3, m\'aximo 5. Clasif\'iquelas en estructurales, operativas, sociales o situacionales.
\item \textbf{Consecuencias} (copa): m\'inimo 3, m\'aximo 5. Ord\'enelas de las m\'as inmediatas a las m\'as amplias.
\end{itemize}
\end{cajaPaso}

\begin{cajaPaso}[\textbf{Componente 2 --- An\'alisis costo-beneficio}]
Eval\'ue la viabilidad de una alternativa de soluci\'on al problema identificado. Puede elegir entre:
\begin{itemize}[leftmargin=1.5em]
\item \textbf{An\'alisis num\'erico}: tabla de costos (iniciales y recurrentes) y beneficios (monetarios y no monetarios), con c\'alculo de \textbf{B/C}, \textbf{VPN} o \textbf{TIR}.
\item \textbf{An\'alisis anal\'itico}: argumentaci\'on en texto explicando por qu\'e los beneficios superan los costos, usando conceptos de la unidad y referenciando literatura.
\end{itemize}
\end{cajaPaso}

\begin{cajaInfo}[\textbf{Formato de entrega}]
Documento PDF de \textbf{2 a 5 p\'aginas}, con portada que incluya nombre completo, c\'edula y secci\'on. Fuente Arial o Times 11--12 pt. Sub\'alo a la plataforma Moodle antes de la fecha de cierre.
\end{cajaInfo}

\newpage

% ==================== 3. EJERCICIO RESUELTO ====================
\section{Ejercicio resuelto de programaci\'on lineal (ejemplo gu\'ia)}

Este ejemplo no es parte de su entrega, pero le servir\'a de \textbf{modelo} para construir un an\'alisis similar en el desarrollo de su actividad. Lo vimos en la clase con un caso del Comando Suroccidente (Cali). Aqu\'i trabajamos un \textbf{caso distinto} para que pueda seguir la estructura.

\subsection{Caso: Operativo de Fin de Semana en Pereira}

\begin{cajaInfo}[\textbf{Situaci\'on planteada}]
El Comandante del Distrito de Pereira debe planificar el \textbf{operativo del fin de semana}. Tiene disponibles \textbf{15 motos} y un presupuesto diario de \textbf{\$2.400.000} para asignarlas a tres sectores:

\begin{itemize}[leftmargin=1.5em]
\item \textbf{Sector 1 --- Centro}: alta afluencia comercial, reportes de hurto a persona concentrados all\'i. Prioridad alta (peso 3), costo de operaci\'on por moto-d\'ia \$160.000, se requieren al menos 5 motos.
\item \textbf{Sector 2 --- Cerritos}: sector residencial de crecimiento, riesgo moderado. Prioridad media (peso 2), costo \$150.000, m\'inimo 3 motos.
\item \textbf{Sector 3 --- Circunvalar}: corredor vial importante, delitos de bajo impacto. Prioridad baja (peso 1), costo \$140.000, m\'inimo 2 motos.
\end{itemize}

\textbf{Pregunta:} \textit{\textquestiondown Cu\'antas motos asignar a cada sector para \textbf{maximizar la cobertura operativa ponderada}, respetando presupuesto y m\'inimos?}
\end{cajaInfo}

\subsection{Datos organizados}

\begin{center}
\begin{tabular}{@{}lccc@{}}
\toprule
\textbf{Sector} & \textbf{Peso prioridad} & \textbf{Costo moto-d\'ia} & \textbf{M\'inimo operativo} \\
\midrule
S1 Centro & 3 & \$160.000 & $\geq 5$ motos \\
S2 Cerritos & 2 & \$150.000 & $\geq 3$ motos \\
S3 Circunvalar & 1 & \$140.000 & $\geq 2$ motos \\
\bottomrule
\end{tabular}
\end{center}

\subsection{Desarrollo paso a paso}

\begin{cajaPaso}[\textbf{Paso 1 --- Definir las variables de decisi\'on}]
Nombre lo que usted puede elegir:
\begin{align*}
x_1 &= \text{n\'umero de motos asignadas al Centro}\\
x_2 &= \text{n\'umero de motos asignadas a Cerritos}\\
x_3 &= \text{n\'umero de motos asignadas al Circunvalar}
\end{align*}
Las tres variables deben ser \textbf{enteros no negativos} (no existen medias motos).
\end{cajaPaso}

\begin{cajaPaso}[\textbf{Paso 2 --- Plantear la funci\'on objetivo}]
Queremos \textbf{maximizar la cobertura ponderada}. Cada sector aporta seg\'un su peso:
\begin{equation*}
\boxed{\; Z = 3\,x_1 + 2\,x_2 + 1\,x_3 \;}
\end{equation*}
Una moto en el Centro aporta \textit{tres veces} lo que una moto en el Circunvalar, porque all\'i el riesgo es mayor.
\end{cajaPaso}

\begin{cajaPaso}[\textbf{Paso 3 --- Escribir las restricciones}]
Cada l\'imite operativo se traduce en una desigualdad. Se escriben \textbf{una debajo de otra}, nunca separadas por comas:

\textbf{Restricci\'on 1 --- Total de motos disponibles:}
\begin{equation*}
x_1 + x_2 + x_3 \leq 15
\end{equation*}

\textbf{Restricci\'on 2 --- Presupuesto diario:}
\begin{equation*}
160\,000\,x_1 + 150\,000\,x_2 + 140\,000\,x_3 \leq 2\,400\,000
\end{equation*}

\textbf{Restricciones 3--5 --- Demandas m\'inimas operativas:}
\begin{align*}
x_1 &\geq 5 \quad\text{(m\'inimo en Centro)}\\
x_2 &\geq 3 \quad\text{(m\'inimo en Cerritos)}\\
x_3 &\geq 2 \quad\text{(m\'inimo en Circunvalar)}
\end{align*}

\textbf{Restricci\'on 6 --- No-negatividad e integralidad:}
\begin{equation*}
x_1, x_2, x_3 \in \mathbb{Z}^{+} \cup \{0\}
\end{equation*}
\end{cajaPaso}

\begin{cajaPaso}[\textbf{Paso 4 --- Consolidar el modelo}]
Se re\'une todo en un sistema, que es lo que se introduce en Excel Solver:
\begin{equation*}
\begin{aligned}
\text{Max}\quad Z &= 3\,x_1 + 2\,x_2 + 1\,x_3\\
\text{s.\,a.}\quad
x_1 + x_2 + x_3 &\leq 15\\
160\,000\,x_1 + 150\,000\,x_2 + 140\,000\,x_3 &\leq 2\,400\,000\\
x_1 &\geq 5\\
x_2 &\geq 3\\
x_3 &\geq 2\\
x_1, x_2, x_3 &\in \mathbb{Z}_{\geq 0}
\end{aligned}
\end{equation*}
\end{cajaPaso}

\begin{cajaPaso}[\textbf{Paso 5 --- Resolver con Excel Solver}]
En Excel: celdas para las variables, una celda con la funci\'on objetivo (uso de \texttt{=SUMAPRODUCTO}) y una por cada restricci\'on. Luego men\'u \textbf{Datos} $\rightarrow$ \textbf{Solver} $\rightarrow$ m\'etodo \textbf{Simplex LP} $\rightarrow$ \textbf{Resolver}. La soluci\'on \'optima es:
\begin{align*}
x_1 &= 9 \quad\text{motos en Centro}\\
x_2 &= 3 \quad\text{motos en Cerritos}\\
x_3 &= 3 \quad\text{motos en Circunvalar}\\
Z &= 3(9) + 2(3) + 1(3) = 36 \quad\text{(cobertura ponderada m\'axima)}
\end{align*}
\end{cajaPaso}

\begin{cajaPaso}[\textbf{Paso 6 --- Verificar y traducir a decisi\'on operativa}]
Se verifica que \textit{todas} las restricciones se cumplen:
\begin{align*}
9 + 3 + 3 &= 15 \leq 15 \quad\checkmark \quad\text{(total de motos)}\\
160\,000(9) + 150\,000(3) + 140\,000(3) &= 2\,310\,000 \leq 2\,400\,000 \quad\checkmark \quad\text{(margen \$90.000)}\\
9 \geq 5,\; 3 \geq 3,\; 3 \geq 2 &\quad\checkmark \quad\text{(m\'inimos)}
\end{align*}
\begin{cajaExito}
\textbf{Decisi\'on operativa:} asignar \textbf{9 motos al Centro, 3 a Cerritos y 3 al Circunvalar}. Queda un margen de \$90.000 para horas extra o imprevistos del turno. Si cambia la prioridad o el presupuesto, solo se actualizan los n\'umeros y se resuelve de nuevo; la estructura del modelo es la misma.
\end{cajaExito}
\end{cajaPaso}

\begin{cajaAlerta}
\textbf{Advertencia metodol\'ogica:} los pesos (3, 2, 1) y los m\'inimos operativos son decisiones del comandante, no son valores autom\'aticos. Deben estar \textbf{justificados} con datos (tasas delictivas, tiempo de respuesta hist\'orico, cifras del SIEDCO). Un modelo sin justificaci\'on de los par\'ametros no es defendible ante el mando.
\end{cajaAlerta}

\newpage

% ==================== 4. PLANTILLA ARBOL ====================
\section{Plantilla de \'arbol de problemas}

A continuaci\'on se muestra un \'arbol de ejemplo (situaci\'on hipot\'etica: \textit{``Alto consumo de combustible en el parque automotor del distrito''}) y luego la plantilla en blanco para que usted la reemplace con su propio caso.

\subsection{Ejemplo completo}

\begin{center}
\begin{tikzpicture}[
  node distance=0.8cm,
  consecuencia/.style={rectangle, rounded corners, fill=consecuencia!20, draw=consecuencia, thick, text width=3.2cm, align=center, font=\footnotesize, minimum height=0.9cm},
  tronco/.style={rectangle, rounded corners, fill=policiaAzul, draw=policiaDorado, very thick, text width=10cm, align=center, font=\small\bfseries\color{white}, minimum height=1.4cm},
  causa/.style={rectangle, rounded corners, fill=causa!20, draw=causa, thick, text width=3.2cm, align=center, font=\footnotesize, minimum height=0.9cm},
  linea/.style={-, thick}
]

% Consecuencias
\node[consecuencia] (c1) at (-5,4) {Mayor gasto\\presupuestal};
\node[consecuencia] (c2) at (0,4) {Reducci\'on de\\horas de patrullaje};
\node[consecuencia] (c3) at (5,4) {Desgaste acelerado\\del parque};

% Etiqueta lateral consecuencias
\node[rotate=90, anchor=east, font=\footnotesize\bfseries\color{consecuencia}] at (-8,4) {CONSECUENCIAS};

% Problema central
\node[tronco] at (0,1.5) {Alto consumo de combustible en el parque automotor del Distrito (incremento del 18\% en el \'ultimo trimestre)};

% Etiqueta lateral problema
\node[rotate=90, anchor=east, font=\footnotesize\bfseries\color{policiaAzul}] at (-8,1.5) {PROBLEMA};

% Causas
\node[causa] (a1) at (-5,-1) {Rutas no planificadas};
\node[causa] (a2) at (0,-1) {Falta de\\mantenimiento preventivo};
\node[causa] (a3) at (5,-1) {Conducci\'on\\inadecuada};

% Etiqueta lateral causas
\node[rotate=90, anchor=east, font=\footnotesize\bfseries\color{causa}] at (-8,-1) {CAUSAS};

% Conexiones tronco-consecuencias
\draw[linea, consecuencia] (c1.south) -- ++(0,-0.3) -| (-3,2.2);
\draw[linea, consecuencia] (c2.south) -- (0,2.2);
\draw[linea, consecuencia] (c3.south) -- ++(0,-0.3) -| (3,2.2);

% Conexiones tronco-causas
\draw[linea, causa] (-3,0.8) |- (a1.north);
\draw[linea, causa] (0,0.8) -- (a2.north);
\draw[linea, causa] (3,0.8) |- (a3.north);

\end{tikzpicture}
\end{center}

\subsection{Plantilla en blanco}

Use el siguiente esqueleto para armar el \'arbol con su caso. Pu\'edalo dibujar a mano, hacerlo en PowerPoint, Draw.io o directamente en LaTeX si conoce TikZ.

\vspace{0.5cm}

\begin{center}
\begin{tikzpicture}[
  node distance=0.8cm,
  hueco/.style={rectangle, rounded corners, fill=grisClaro, draw=grisBorde, thick, text width=3.2cm, align=center, font=\footnotesize\color{gray}, minimum height=1.1cm},
  troncoBlanco/.style={rectangle, rounded corners, fill=white, draw=policiaAzul, very thick, text width=10cm, align=center, font=\small\color{gray}, minimum height=1.6cm}
]
\node[hueco] (c1) at (-5,4) {[Consecuencia 1]};
\node[hueco] (c2) at (0,4) {[Consecuencia 2]};
\node[hueco] (c3) at (5,4) {[Consecuencia 3]};

\node[troncoBlanco] at (0,1.5) {[Escriba aqu\'i el problema central, \\una frase clara, medible y ubicada en el tiempo y el territorio]};

\node[hueco] (a1) at (-5,-1) {[Causa 1]};
\node[hueco] (a2) at (0,-1) {[Causa 2]};
\node[hueco] (a3) at (5,-1) {[Causa 3]};

\draw[-, thick, gray] (c1.south) -- ++(0,-0.3) -| (-3,2.3);
\draw[-, thick, gray] (c2.south) -- (0,2.3);
\draw[-, thick, gray] (c3.south) -- ++(0,-0.3) -| (3,2.3);
\draw[-, thick, gray] (-3,0.7) |- (a1.north);
\draw[-, thick, gray] (0,0.7) -- (a2.north);
\draw[-, thick, gray] (3,0.7) |- (a3.north);
\end{tikzpicture}
\end{center}

\begin{cajaInfo}[\textbf{Consejo}]
Empiece por el \textbf{problema central}. Una frase bien formulada responde tres preguntas: \textit{qu\'e ocurre}, \textit{d\'onde} y \textit{cu\'ando}. Evite frases como ``hay inseguridad''; prefiera ``incremento del 22\% de hurtos a persona en el cuadrante 3 entre enero y marzo de 2026''.
\end{cajaInfo}

\newpage

% ==================== 5. COSTO-BENEFICIO ====================
\section{Plantilla de an\'alisis costo-beneficio}

\subsection{Fundamentos te\'oricos r\'apidos}

\begin{center}
\begin{tabularx}{\textwidth}{@{}lXl@{}}
\toprule
\textbf{Indicador} & \textbf{\textquestiondown Qu\'e significa?} & \textbf{F\'ormula} \\
\midrule
\textbf{B/C} & Raz\'on beneficios/costos. Dice cu\'antos pesos de beneficio produce cada peso invertido. Es la forma m\'as intuitiva de presentar al mando. & $\dfrac{\sum B}{\sum C}$ \\[8pt]
\textbf{VPN} & Valor Presente Neto. Suma de flujos netos descontados al presente. Dice cu\'anto vale la inversi\'on hoy. & $\displaystyle\sum_{t=0}^{n} \dfrac{B_t - C_t}{(1+i)^t}$ \\[10pt]
\textbf{TIR} & Tasa Interna de Retorno. Rentabilidad anual impl\'icita del proyecto. Es la tasa que hace que el VPN sea cero. & $i^{*}$ tal que $VPN(i^{*})=0$ \\
\bottomrule
\end{tabularx}
\end{center}

\textbf{Criterios de decisi\'on:}
\begin{itemize}[leftmargin=1.5em]
\item B/C $> 1$, VPN $> 0$, TIR $>$ tasa referencia $\;\Rightarrow\;$ la inversi\'on es viable.
\item B/C $= 1$, VPN $= 0$, TIR $=$ tasa referencia $\;\Rightarrow\;$ la inversi\'on es indiferente.
\item En caso contrario, no es viable.
\end{itemize}

\subsection{Ejemplo num\'erico}

\textbf{Caso:} compra de 5 radios de comunicaci\'on nuevos versus reparaci\'on de los 5 actuales.

\begin{center}
\begin{tabularx}{\textwidth}{@{}lXX@{}}
\toprule
\textbf{Concepto} & \textbf{Opci\'on A: comprar nuevos} & \textbf{Opci\'on B: reparar los actuales}\\
\midrule
Inversi\'on inicial & \$12.000.000 & \$3.000.000\\
Costo operativo anual & \$600.000 & \$2.200.000 (mantenimientos recurrentes)\\
Beneficios anuales & \$5.500.000 (menos fallas, mejor cobertura) & \$3.000.000\\
Horizonte & 5 a\~nos & 2 a\~nos (vida \'util restante)\\
Tasa descuento & 10\% & 10\%\\
\midrule
\textbf{VPN} & \textbf{+\$6.562.000} $\checkmark$ & \textbf{+\$388.000}\\
\textbf{B/C} & \textbf{1,54} & \textbf{1,06}\\
\bottomrule
\end{tabularx}
\end{center}

\begin{cajaExito}
\textbf{Decisi\'on:} la Opci\'on A domina en rentabilidad, aunque la inversi\'on inicial es mayor. Si el distrito cuenta con el capital para el desembolso inicial, es la opci\'on m\'as eficiente a mediano plazo.
\end{cajaExito}

\subsection{Plantilla en blanco para el estudiante}

Complete con su propio caso. Puede hacer un an\'alisis \textbf{num\'erico} (llenando los valores) o \textbf{anal\'itico} (reemplazando los n\'umeros por argumentos textuales).

\begin{center}
\begin{tabularx}{\textwidth}{@{}lXX@{}}
\toprule
\textbf{Concepto} & \textbf{Alternativa A} & \textbf{Alternativa B / Status quo}\\
\midrule
Descripci\'on breve & \rule{0pt}{1.2em}\hrulefill & \rule{0pt}{1.2em}\hrulefill\\[6pt]
Inversi\'on inicial & \hrulefill & \hrulefill\\[6pt]
Costos recurrentes & \hrulefill & \hrulefill\\[6pt]
Beneficios (monetarios) & \hrulefill & \hrulefill\\[6pt]
Beneficios (no monetarios) & \hrulefill & \hrulefill\\[6pt]
Horizonte (a\~nos) & \hrulefill & \hrulefill\\[6pt]
Tasa descuento & \hrulefill & \hrulefill\\[6pt]
\textbf{VPN / B/C / TIR} & \hrulefill & \hrulefill\\[6pt]
\textbf{Decisi\'on justificada} & \multicolumn{2}{p{0.6\textwidth}}{\rule{0pt}{1.2em}\hrulefill}\\[6pt]
\bottomrule
\end{tabularx}
\end{center}

\begin{cajaInfo}[\textbf{Versi\'on anal\'itica sin n\'umeros}]
Si no tiene cifras confiables, argumente en texto. Un an\'alisis anal\'itico v\'alido responde estas tres preguntas:
\begin{enumerate}[leftmargin=1.5em]
\item \textbf{\textquestiondown Qu\'e costos implica la alternativa?} (iniciales, recurrentes, ocultos, riesgos).
\item \textbf{\textquestiondown Qu\'e beneficios genera?} (monetarios, operativos, humanos, institucionales).
\item \textbf{\textquestiondown Por qu\'e los beneficios superan los costos?} (argumentos sustentados en la literatura o en experiencia documentada).
\end{enumerate}
Extensi\'on sugerida: 300--500 palabras.
\end{cajaInfo}

\newpage

% ==================== 6. RUBRICA ====================
\section{R\'ubrica de evaluaci\'on}

La entrega se eval\'ua sobre \textbf{5,00 puntos} con los siguientes criterios:

\begin{center}
\footnotesize
\begin{tabularx}{\textwidth}{@{}p{3.4cm}cXXXX@{}}
\toprule
\textbf{Criterio} & \textbf{Puntos} & \textbf{Excelente (1,0 / 0,8)} & \textbf{B\'asico (0,6 / 0,3)} & \textbf{No entrega (0)}\\
\midrule
1. An\'alisis del problema y \'arbol
& 1,0
& Identifica claramente el problema central, causas y consecuencias con argumentaci\'on s\'olida
& Identifica el problema pero con an\'alisis superficial o incompleto
& No entrega la actividad\\[6pt]
2. Formulaci\'on matem\'atica (si aplica) o descriptiva clara
& 1,0
& Plantea correctamente la alternativa con c\'alculos o argumentos bien fundados
& La formulaci\'on tiene errores conceptuales o matem\'aticos
& No entrega la actividad\\[6pt]
3. An\'alisis costo-beneficio
& 1,0
& Aplica correctamente el an\'alisis con datos o argumentos detallados
& Presenta un an\'alisis superficial o con fallos metodol\'ogicos
& No entrega la actividad\\[6pt]
4. Conclusiones y justificaci\'on de la decisi\'on
& 1,0
& Justifica claramente la viabilidad de la alternativa
& Justificaci\'on vaga o sin conexi\'on con los datos
& No entrega la actividad\\[6pt]
5. Presentaci\'on y estructura del informe
& 1,0
& Documento claro, bien estructurado, con gr\'aficos organizados
& Estructura adecuada con detalles mejorables
& No entrega la actividad\\
\midrule
\textbf{Total} & \textbf{5,0} & & & \\
\bottomrule
\end{tabularx}
\end{center}

\begin{cajaAlerta}
\textbf{Importante:} la pol\'itica ESJIM fija un piso de \textbf{4,10} y un techo de \textbf{5,00} para las entregas evaluadas con IA. El docente conserva la potestad de ajuste seg\'un el contenido real del trabajo.
\end{cajaAlerta}

% ==================== 7. CRONOGRAMA ====================
\section{Cronograma sugerido}

\begin{center}
\begin{tabular}{@{}lp{9cm}@{}}
\toprule
\textbf{Semana} & \textbf{Actividad recomendada}\\
\midrule
Semana 1 & Revisar la clase grabada y leer la bibliograf\'ia de la unidad.\\
Semana 2 & Elegir la situaci\'on cotidiana del servicio que quiere analizar y hacer un borrador del \'arbol de problemas.\\
Semana 3 & Construir las tablas de costos y beneficios; calcular B/C o VPN, o redactar el an\'alisis anal\'itico.\\
Semana 4 & Redactar el informe final (2--5 p\'aginas), revisar con un compa\~nero y subir a Moodle antes del 14 de mayo, 23:59.\\
\bottomrule
\end{tabular}
\end{center}

% ==================== 8. RECURSOS ====================
\section{Recursos y bibliograf\'ia}

\subsection{Bibliograf\'ia b\'asica (obligatoria)}
\begin{itemize}[leftmargin=1.5em]
\item Pont\'on Cevallos, D. (2020). El aporte de Edwin Sutherland al an\'alisis del crimen econ\'omico global. \textit{URVIO: Revista Latinoamericana de Estudios de Seguridad}, 27, 112--124.
\item Rohatgi, V. K., \& Saleh, A. K. E. (2000). \textit{An Introduction to Probability and Statistics}. John Wiley \& Sons.
\item Villarreal Satama, F. L., Montenegro G\'alvez, D. I., \& N\'u\~nez Ribadeneira, J. E. (2021). Optimizaci\'on Matem\'atica como Herramienta para la toma de decisiones en la Empresa. \textit{Ingenio}, 4(1), 40--60.
\end{itemize}

\subsection{Bibliograf\'ia complementaria}
\begin{itemize}[leftmargin=1.5em]
\item Alvino Aldave, V. M., et al. (2024). Anticipando el Crimen: Modelos Matem\'aticos para la Proyecci\'on del Comportamiento Delincuencial en la Provincia de Huaura. \textit{TecnoHumanismo}, 4(3), 1--169.
\item L\'opez, G. J., Coronel, E. L., \& R\'ios, M. E. (2012). Simulador de programaci\'on lineal 2D y 3D como recurso did\'actico. \textit{Revista Iberoamericana de Ciencias y Tecnolog\'ia}, 3(2), 123--135.
\end{itemize}

\subsection{Herramientas}
\begin{itemize}[leftmargin=1.5em]
\item \textbf{Microsoft Excel} con complemento Solver (\textit{Datos} $\rightarrow$ \textit{Solver} $\rightarrow$ \textit{Simplex LP}).
\item \textbf{Khan Academy}: secciones de probabilidad y optimizaci\'on.
\item \textbf{GeoGebra} para representaci\'on gr\'afica de problemas de dos variables.
\end{itemize}

\subsection{Material de apoyo del docente}
\begin{itemize}[leftmargin=1.5em]
\item Presentaci\'on HTML: \url{https://jotaj1991.github.io/recursos-esjim/u2_plan_clase/plan_clase_u2.html}
\item Plantilla Excel con el caso de Cali resuelto: \texttt{plantilla\_PL\_patrullas.xlsx}
\item Bases de datos institucionales: BD1 Incidentes Bogot\'a, BD2 Gesti\'on Cuadrantes Colombia, BD3 Perfil Criminal y Reincidencia.
\end{itemize}

\vspace{1cm}

\begin{center}
\color{policiaAzul}\rule{0.6\textwidth}{0.6pt}\\[8pt]
\footnotesize\textit{Cualquier duda sobre el desarrollo de la actividad, consulte al docente por el canal oficial de la plataforma Moodle.}
\end{center}

\end{document}
